% MT033 / explanation — canonical editable content.
把图像表示为特征向量\(x\)，标签\(c\)是题目训练数据已经给定的类别。


\textbf{生成式路线：分别认识每类的特征分布。} 若用简单高斯模型，训练会得到每类的平均特征、各特征的变化范围及类别比例。预测新\(x\)时，计算它在每类模型下有多常见，再乘该类先验，通过Bayes公式比较。学出的均值和方差描述题中类别的数据特点。


\textbf{判别式路线：直接学习区分类别的规则。} 例如LR学习各特征的权重\(w\)和偏置\(b\)，先算分数\(w^Tx+b\)，再用sigmoid函数把分数转成0到1之间的类别概率。某个特征的权重为正，表示在其他特征固定时，增大它会提高模型预测为类1的概率；不要求同时学出该类所有图像特征怎样分布。


两者的取舍来自假设。简单生成式模型对分布作较强限制，假设合适时可能用较少带标签样本就学得稳定；假设不合适时也会产生系统偏差。判别式把学习集中在分类上，不必完整描述输入分布，但仍需要足够数据来估计判别规则。


解释模型时还要看特征单位：以像素数或归一化数值表示同一特征，权重大小会改变。权重和特征重要性描述模型如何预测，不直接表示现实因果关系。


\textbf{出处：} 2021B*期中 Q7，原卷第3页（10分）


\textbf{答案依据：} 2021B配套解答 / 2021B配套答案副本 Q7，随卷解答；补充样本效率的条件。
